\documentclass[a4paper]{article}
\usepackage[pdftex]{graphicx}
\usepackage[utf8]{inputenc}
\usepackage{enumerate}
\usepackage{amssymb}
\usepackage{tikz}
\usepackage{href-ul}
\hypersetup{
	colorlinks=true,
	linkcolor=blue,
	urlcolor=blue}
\usepackage{geometry}
\geometry{a4paper, top=15mm, left=15mm, right=15mm, bottom=15mm,
	headsep=10mm, footskip=12mm}

\begin{document}
	\begin{tikzpicture}(10,3)
		\draw[ultra thick](2,0) --(10,0) -- (10,1.3) --(2,1.3) -- (2,0);
		\draw[fill=black](2,0)-- (0,.6) -- (2,.6) -- (2,0);
		\node at (6,.6) {\large solution sheet from \href{http://www.okuyakl.com}{www.okuyakl.com}};
	\end{tikzpicture}
	\vspace{0.5 cm}
	
	\noindent{\bf Task 1. a)}\\
	
	\begin{itemize}
		\item $b$ is the starting value, it can be any real number. We have $b=f(0)$.
		\item $a$ is the growth factor. It must be positive and different from 1:
		$a \in \mathbb{R}^+ \setminus \{1\}.$
	\end{itemize}
	
	\noindent{\bf Task 1. b)}\\
	\begin{itemize}
		\item \underline{Case 1. $a < 0$; e.g. B. $a=-1; \quad x={1\over 2}$}
		$$y=b \cdot a^{1\over 2} = b \cdot \sqrt{-1} \qquad \textnormal{no solution}$$
		\item \underline{Case 2. $a = 0; \quad x=-1$ }
		$$y=b \cdot a^{-1} = b \cdot { 1 \over 0} \qquad \textnormal{no solution}$$
		\item \underline{Case 2. $a = 1$: }
		$$y= b \cdot 1^{x} = b \qquad \textnormal{for all } \qquad x \in \mathbb{R}$$
	\end{itemize}
	
	\noindent{\bf Task 1. c)}\\
	\begin{enumerate}[(i)]
		\item $$f(x)= 2\cdot \left({1\over 2}\right)^x$$
		The function is decreasing because $a<1$ and the exponent is positive. We have $f(0)=2 \not = 1$.
		\item $$f(x)= 2 \cdot 2^x $$
		The function is increasing because $a>1$ and the exponent is positive. We have $f(0)=2$.
	\end{enumerate}
	
	\noindent
	\fbox{
		\begin{minipage}{0.3\textwidth}
			\noindent{\bf Task 2. a)}\\
			Exponent comparison:
			$$
			\renewcommand{\arraystretch}{2}
			\begin{array}{rcll}
				4^{2x-5} &=& 4^2 \\
				2x-5 &=& 2 \\
				x &=& 3.5
			\end{array}
			$$
		\end{minipage}
	}
	\hfill
	\fbox{
		\begin{minipage}{0.3\textwidth}
			\noindent{\bf Task 2. b)}\\
			$$
			\renewcommand{\arraystretch}{2}
			\begin{array}{rcll}
				3^x &=& 27 \\
				3^x &=& 3^3 \\
				x &=& 3
			\end{array}
			$$
		\end{minipage}
	}
	\hfill
	\fbox{
		\begin{minipage}{0.3\textwidth}
			\noindent{\bf Task 2. c)}\\
			$$
			\renewcommand{\arraystretch}{2}
			\begin{array}{rcll}
				2^{x^2-x-7} &=& 2^{-5} \\
				x^2 -x -7 &=& -5 \\
				x^2 -x -2 &=& 0 \\
				x_1 &=& 2 \\
				x_2 &=& -1
			\end{array}
			$$
		\end{minipage}
	}
	
	\noindent{\bf Task 3.}\\
	We plug in the points and get two equations:
	$$
	\renewcommand{\arraystretch}{2}
	\begin{array}{rccll}
		I & 36 &=& b \cdot a^2 \\
		II & {32 \over 3} &=& b \cdot a^5 \\
		\hline
		II : I & {8 \over 27} &=& a^3 \\
		& a &=& {2 \over 3} \\
		& b &=& {36 \cdot 9 \over 4} &= 81\\
		\Rightarrow & f(x)&=& 81 \cdot \left({2 \over 3} \right)^x
	\end{array}
	$$
	
	\vspace{0.5 cm}
	
	\noindent{\bf Task 4.}\\
	\renewcommand{\arraystretch}{2}
	\begin{tabular}{|p{1.5 cm}|p{1.5 cm}|p{1.5 cm}|p{1.5 cm}|p{1.5 cm}|p{1.5 cm}|p{1.5 cm}|p{1.5 cm}|p{1.5 cm}|}
		\hline
		$a$ & 1.02 & 1.1 & 1.01 & 1.35 & 1.07 & 1.05 & 2 & 1.25 \\
		\hline
		$p$ & 2\% & 10\% & 1\% & 35\% & 7\% & 5\% & 100\%& 25\% \\
		\hline
		$T$ & 35 & 7.3 & 70 & 2.3 & 10 & 14 & 1 & 3.1 \\
		\hline
	\end{tabular}
	\vspace{0.5 cm}
	
	\noindent{\bf Task 5.}\\
	
	\renewcommand{\arraystretch}{2}
	\begin{tabular}{|p{1.5 cm}|p{1.5 cm}|p{1.5 cm}|p{1.5 cm}|p{1.5 cm}|p{1.5 cm}|p{1.5 cm}|p{1.5 cm}|p{1.5 cm}|}
		\hline
		$a$ & 0.97 & 0.92 & 0.93 & 0.8 & 0.5 & 0.965 & 0.995 & 0.86 \\
		\hline
		$p$ & 3\% & 8\% & 7\% & 20\% & 50\%& 3.5\% & 0.5\% & 14\% \\
		\hline
		$T$ & 23 & 8.3 & 10 & 3.1 & 1 & 20 & 138 & 4.6 \\
		\hline
	\end{tabular}
	
	\begin{center}
		\includegraphics[width=7 cm]{../../viecher/eendcomic.pdf}
		
		Here you can go back to the \href{https://www.okuyakl.de/math/m10expfuL030/ae030.pdf}{task sheet}
	\end{center}
	
\end{document}